\section{What straightedge and compass cannot construct}
\label{sec:obstructions}
Fix the rational Cartesian frame $(0,0),(1,0),(0,1)$.  A construction
program is a finite list of the operations in
Theorem~\ref{circ:construction-arithmetic}, with nonzero-denominator and
intersection-validity data.  Thus its coordinate recipe uses rational
arithmetic and successive nonnegative square roots.  The number $k$ below
counts these root adjunctions, including any used to measure an output
segment.  A repeated root or a root already in the current field need not
increase its dimension.  The classical construction obstruction is due
to Wantzel; we give the finite algebra needed for its arithmetic use.\footnote{
L. Wantzel, ``Recherches sur les moyens de reconna\^itre si un Probl\`eme
de G\'eom\'etrie peut se r\'esoudre avec la r\`egle et le compas'',
\emph{Journal de Math\'ematiques Pures et Appliqu\'ees}, series 1,
vol.~2 (1837), pp.~366--372;
\url{https://www.numdam.org/item/JMPA_1837_1_2__366_0/}.}

\begin{theorem}[The diagonal of the unit square is irrational]
\label{obs:diagonal}
The unit square has a constructible diagonal of length $\sqrt2$, and no
rational number has square $2$.
\end{theorem}
\begin{proof}
The diagonal vector is $(1,1)$, whose squared length is $2$; its positive
length is therefore $\sqrt2$.  The vertices, and hence the segment, are
constructed by the perpendicular and parallel constructions already proved.
If a rational of square $2$ existed, write it as $m/n$ with integers
$m\ge0$, $n>0$, and $\gcd(m,n)=1$.  Reducing a rational to this form
uses the bounded integer Euclidean algorithm.  The equality would be
$m^2=2n^2$.  Squares have the parity of their bases, since
$(2a)^2$ is even and $(2a+1)^2$ is odd.  Thus $m=2a$, and substitution
gives $n^2=2a^2$, so $n$ is even too.  This contradicts coprimality.
Equivalently, a computed B\'ezout relation $um+vn=1$ would have an even
left side.  The argument is a finite parity calculation, not an infinite
descent or a minimum principle over arbitrary sets.
\end{proof}

\subsection{Finite field bases and the degree obstruction}
A finite number-field record consists of a monic irreducible rational
polynomial $h$ and a selected scalar $\xi\in\K$ with $h(\xi)=0$.
Its elements are coefficient vectors for $1,\xi,\ldots,\xi^{d-1}$,
where $d=\deg h$.  Polynomial remainder computes multiplication.
Extended Euclid computes the inverse of a nonzero vector: irreducibility
makes its representing polynomial coprime to $h$.
The bounded factorization procedure in Lemma~\ref{found:convert} tests
irreducibility and supplies the required factor certificates.

Here is an explicit bounded way to merge two such fields, so that
``find a basis'' does not hide an unbounded search.  For generators
$\alpha,\beta$ with monic irreducible polynomials $p,q$ of degrees $m,n$,
use the rational algebra
\[
 E=\Q[U,V]/(p(U),q(V)),\qquad N=mn,
\]
with its known $N$ monomials as basis.  Let $H_c$ be the characteristic
polynomial of multiplication by $U+cV$.  Its discriminant, as a
polynomial in $c$, has degree at most $M=m^2n(n-1)$, and its coefficient
of $c^M$ is
\[
\operatorname{disc}(p)^n\operatorname{disc}(q)^{m^2}\ne0.
\]
Here the discriminant of a monic degree-$d$ polynomial is
$(-1)^{d(d-1)/2}$ times the Sylvester determinant of it and its derivative,
so it is computed by finite coefficient arithmetic.
Irreducibility in characteristic zero gives
$\gcd(p,p')=\gcd(q,q')=1$, since each derivative is nonzero and has
smaller degree; extended Euclid and the Sylvester determinant identity
therefore certify that both discriminants are nonzero.
To verify this finite coefficient identity, introduce formal roots
$a_i,b_j$.  Evaluation at their pairs has an invertible tensor product
of Vandermonde matrices over the rational-function field, so the
characteristic polynomial is formally $\prod_{i,j}(T-a_i-cb_j)$.
In the product
of squared pairwise differences, pairs with equal $j$ contribute
$\operatorname{disc}(p)^n$, and pairs with different $j$ contribute the
leading factor $\operatorname{disc}(q)^{m^2}c^{m^2n(n-1)}$.
Symmetric coefficient comparison gives the stated identity for $p,q$.
This substitution is injective: products of the elementary symmetric
polynomials have distinct leading monomials, whose exponent vectors are
the successive tail sums of the product exponents.  Thus a coefficient
identity after formal-root substitution is an identity before it;
no choice of actual complex roots is required.  A nonzero polynomial of
degree at most $M$ has at most $M$ roots, by successive division by linear
factors.  Checking $c=0,\ldots,M$ therefore finds one with $H_c$ squarefree.

For this $c$, the powers of $U+cV$ form a basis of $E$.  Indeed, the
matrix with entries
$\operatorname{Tr}(A^{i+j})$, $0\le i,j<N$, where $A$ is its
multiplication matrix, has determinant $\operatorname{disc}(H_c)$.
Trace means the sum of the diagonal entries.  Newton's coefficient
identities identify these traces with the formal power sums, so the
determinant identity is the square of a Vandermonde determinant.
A dependence among the powers
would annihilate this nonsingular matrix, so none exists.  Rational
Gaussian elimination therefore expresses $U,V$ as polynomials in $U+cV$.
Factor $H_c$, and select the unique irreducible factor $h$ vanishing at
$\xi=\alpha+c\beta$, using the exact scalar equality test.  Such a factor
exists because their product vanishes, and two distinct factors cannot
both vanish by B\'ezout.  Reducing the recovered expressions modulo $h$
gives a power-basis record for $\Q(\alpha,\beta)=\Q(\xi)$.
All searches have just been bounded: the $M+1$ choices of $c$, the
factorization enumeration, and elimination on explicitly dimensioned
rational matrices.  Repeated merges are primitive recursion on a supplied
finite generator list.

\begin{theorem}[Quadratic-tower degree obstruction]
\label{obs:degree}
Let
\[
 L_0=\Q,\qquad L_{i+1}=L_i(\beta_i),\qquad
 \beta_i^2=a_i\in L_i\quad(0\le i<k)
\]
be a finite tower of real algebraic fields.  Then
$[L_k:\Q]=2^j$ for some $j\le k$.  For every $x\in L_k$, its
minimal-polynomial degree $e=[\Q(x):\Q]$ divides $2^j$, and hence
divides $2^k$.  These degree statements have a uniform numerical proof
for the finite field-and-basis certificates just described; every finite
straightedge--compass recipe produces such certificates by PR operations.
\end{theorem}
\begin{proof}
First, $L_i(\beta_i)$ is spanned over $L_i$ by $1,\beta_i$.
The relation $\beta_i^2=a_i$ reduces products to this span.  For an
inverse use
\[
 (u+v\beta_i)^{-1}=\frac{u-v\beta_i}{u^2-a_iv^2}
\]
when the denominator is nonzero.  If that denominator is zero and the
original element is nonzero, the zero-product law forces $u-v\beta_i=0$;
unless $v=0$, this says $\beta_i=u/v\in L_i$.  The $v=0$ case is
an ordinary inverse in $L_i$.  Thus the span is a field in every case.

The field records above make membership $\beta_i\in L_i$ decidable by
a rational linear system: express the old basis and $\beta_i$ in the
new power basis and test the rank after adjoining its column.  If
membership holds, the old basis remains a basis.  Otherwise $1,\beta_i$
are linearly independent over $L_i$: a relation $u+v\beta_i=0$ with
$v\ne0$ would express $\beta_i=-u/v\in L_i$, while $v=0$ forces $u=0$.
For an old rational basis $e_1,\ldots,e_D$, the resulting new rational
basis is therefore either the old one or
\[
 e_1,\ldots,e_D,\quad \beta_i e_1,\ldots,\beta_i e_D.
\]
Hence the dimension is unchanged or doubled.  Induction gives
$D_i=2^{j_i}$ with $j_i\le i$, counting only actual doublings.

For completeness, divisibility is also a finite basis calculation.
In a field of rational dimension $D$, the $D+1$ vectors
$1,x,\ldots,x^D$ are dependent.  Elimination for degrees $1,\ldots,D$
finds the first monic relation $\mu(x)=0$, of degree $e$.
If $\mu=fg$ with both factor degrees smaller, the zero-product law
would contradict its minimality.  Thus $\mu$ is irreducible, and
$1,x,\ldots,x^{e-1}$ is a basis of the subfield $\Q(x)$.
Scan the given rational basis of $L_k$, greedily retaining a vector $w$
when it is not in the $\Q(x)$-span of earlier retained vectors.
Each test is rational Gaussian elimination on the columns $x^r w$,
$0\le r<e$.  At most $D$ vectors are scanned.  If $f$ are retained,
the $ef$ products $x^r w$ are a rational basis: spanning follows from
the scan, and any dependence groups into a $\Q(x)$-dependence of the
retained vectors.  Thus $D=ef$, proving $e\mid2^j$.
Multiplication by $2^{k-j}$ proves $e\mid2^k$.

The invariant $D_i=2^{j_i}$, $j_i\le i$, together with the verified
coefficient-vector and rank certificates, is a quantifier-free statement
about a finite list code.  Its induction is legitimate in \ipra.
The bounded field-merge construction supplies the certificates at each
root instruction, while rational operations stay in the current field.
This proves the asserted uniformity without quantifying over arbitrary
vector spaces or choosing an unspecified basis.
\end{proof}

\paragraph{A finite irreducibility test used three times.}
A cubic over $\Q$ is reducible exactly when it has a rational root: a
nontrivial factorization of degree three has a linear factor, and a root
gives such a factor by division.  A rational root $u/v$ in lowest terms
of a monic integer polynomial has $v=1$, since the cleared equation
gives $v\mid u^d$.  Raising a B\'ezout relation for $u,v$ to the $d$th
power then gives $v\mid1$, hence $v=1$.  The resulting integer
root divides the constant coefficient.  Thus for nonzero constant term
$c$ only the signed divisors in $[-|c|,|c|]$ need checking.  These are
bounded arithmetic tests.  Finally $3$ never divides $2^j$: its residue
alternates between $1$ and $2$, a quantifier-free induction on $j$.

\begin{theorem}[Doubling the unit cube is impossible]
\label{obs:cube}
No finite straightedge--compass construction in the rational initial frame
produces the edge of a cube of twice the unit cube's volume.
\end{theorem}
\begin{proof}
Such a positive edge length $x$ would satisfy $x^3=2$.
The only possible rational roots of $T^3-2$ are $\pm1,\pm2$, and
none is a root.  The polynomial is therefore irreducible of degree three.
Theorem~\ref{obs:degree} would make $3$ divide a power of two, a
contradiction.  The desired length does exist in $\K$, by the polynomial
intermediate value theorem between $1$ and $2$; the obstruction concerns
the specified finite geometric operations.
\end{proof}

\begin{theorem}[A $60$-degree angle cannot be trisected in general]
\label{obs:trisection}
Starting with the rational frame and its constructible $60$-degree ray,
no finite straightedge--compass program constructs an internal trisector
making a $20$-degree angle with the base ray.  Consequently there is no
such program that trisects every supplied angle.
\end{theorem}
\begin{proof}
The given angle pair is $(1/2,\sqrt3/2)$; adjoining its coordinates costs
at most one quadratic extension.  Normalize a proposed trisector to a
unit direction $(c,s)$, using one further nonnegative square root if
necessary.  The algebraic angle multiplication of
Section~\ref{sec:triangles} gives
\[
 (c+is)^3=\frac12+i\frac{\sqrt3}{2},\qquad
 c^2+s^2=1,
\]
where $i$ is only notation for the pair product.  Comparing real parts
gives $4c^3-3c=1/2$.  Hence $t=2c$ satisfies
$t^3-3t-1=0$.  Its only possible rational roots are $1,-1$, where the
polynomial has values $-3,1$.  Its degree is therefore three, contradicting
Theorem~\ref{obs:degree}.

The angle itself exists in the coded field: the cubic has a root
$1<t<2$, because its values at the endpoints have opposite signs.
Set $c=t/2$, $s=\sqrt{1-c^2}>0$.  The cube of this unit pair has real
part $1/2$ and positive imaginary part $s(t^2-1)$, hence is the stated
$60$-degree pair.  Its first two successive products are in the upper
semicircle, so the angle-addition convention has no half-turn carry.
This supplies the intended internal third.  The impossibility assertion
is thus the failure of one concrete constructible input, which suffices
to refute a general trisection procedure.
\end{proof}

\begin{theorem}[A regular heptagon is not constructible]
\label{obs:heptagon}
No finite straightedge--compass construction in the rational frame
inscribes a regular heptagon in the unit circle.  More generally it cannot
produce a nondegenerate regular heptagon in any location or scale.
\end{theorem}
\begin{proof}
After setting one vertex at $(1,0)$, let the rotation to its next vertex
be the unit pair $z$.  Regularity here has the equal-successive-rotation,
convex-boundary-order meaning of Theorem~\ref{poly:regular}.  Consequently
$z^7=1$ and $z\ne1$.
Nonzero pairs are cancellable, since their squared norms multiply and
are positive.  Therefore
\[
 1+z+z^2+z^3+z^4+z^5+z^6=0.
\]
Multiply by $z^{-3}$ and put $t=z+z^{-1}$, which is twice the real
coordinate of $z$.  The identities
$z^2+z^{-2}=t^2-2$ and $z^3+z^{-3}=t^3-3t$ give
\[
 t^3+t^2-2t-1=0.
\]
The values at the only possible rational roots $1,-1$ are $-1,1$.
Thus this cubic is irreducible, and again degree three contradicts
Theorem~\ref{obs:degree}.
For a heptagon in an arbitrary position, compute its circumcenter from
three vertices, translate it to zero, divide by its positive radius, and
rotate its first unit vertex to $(1,0)$.  The circumcenter formulas,
square root for the radius, and unit-frame rotation use only permitted
field operations and square roots.  Consecutive vertices then give the
unit rotation used above.  This reduces every nondegenerate regular
heptagon to the already impossible case.
\end{proof}

\paragraph{Scope and formalization.}
These are arithmetic obstructions to finite constructions, with the
initial data specified.  A Lean implementation should encode the
irreducible-polynomial records, rational matrices, field embeddings,
and the two possible basis updates; then prove the dimension invariant
by induction on the construction list.  The three displayed cubic
identities and their finite rational-root checks are the final short
steps.  The present package does not claim a compiled implementation of
this additional number-field certificate layer.  Circle squaring would
require the separate transcendence argument for $\pi$ and is not claimed
as a consequence of these cubic calculations.
